\documentclass[11pt,a4paper]{article}

\usepackage[utf8]{inputenc}
\usepackage[T1]{fontenc}
\usepackage{lmodern}
\usepackage[margin=2.2cm]{geometry}
\usepackage{xcolor}
\usepackage{titlesec}
\usepackage{enumitem}
\usepackage{parskip}
\usepackage[hidelinks]{hyperref}
\usepackage{fancyhdr}

% --- styling ---------------------------------------------------------------
\definecolor{accent}{RGB}{20,86,140}
\definecolor{muted}{RGB}{90,90,90}

\titleformat{\section}{\large\bfseries\color{accent}}{\thesection.}{0.5em}{}
\titlespacing*{\section}{0pt}{1.4ex plus .2ex}{0.6ex}

\setlist[enumerate]{itemsep=1pt,topsep=2pt,parsep=0pt}
\setlist[itemize]{itemsep=1pt,topsep=2pt,parsep=0pt}

\pagestyle{fancy}
\fancyhf{}
\renewcommand{\headrulewidth}{0pt}
\fancyfoot[L]{\footnotesize\color{muted}Strategic Technology Foresight --- Sprint Report}
\fancyfoot[R]{\footnotesize\color{muted}\thepage}

\hypersetup{pdftitle={Sprint Report --- Strategic Technology Foresight},
            pdfauthor={Paul Keck}}

\begin{document}

% --- title block -----------------------------------------------------------
\begin{center}
  {\LARGE\bfseries\color{accent} Sprint Report}\\[3pt]
  {\large Strategic Technology Foresight}
\end{center}
\vspace{2pt}

\noindent\textbf{Project:} AI-Driven Scaling of Strategic Technology Foresight
(Digital Innovation Lab, partner: Rohde \& Schwarz)\\[2pt]
\textbf{Author:} Paul Keck \quad\textbullet\quad \textbf{Period:} 20--23 June 2026
\quad\textbullet\quad \textbf{Audience:} Scrum Master (non-technical summary)

\vspace{-1ex}
\noindent\textcolor{accent}{\rule{\linewidth}{0.7pt}}

% --- body ------------------------------------------------------------------
\section{In a nutshell}
This sprint had two big outcomes. First, we found and fixed a hidden quality problem: the
futures our AI was generating \emph{looked} varied but were, on closer inspection, close to
random noise --- and we can now prove, with a measurable test, that the fix produces real
substance. Second --- and more important --- we clarified with the product owner what the
project is \emph{really} supposed to deliver and rebuilt the tool to match it: a
\textbf{reusable engine that works for any topic}, not just the one demo topic. We proved this
works by plugging in a completely different subject area and getting sensible results with
\textbf{zero reprogramming}. Everything is covered by 109 automated tests.

\section{What the project is about}
Companies like our partner Rohde \& Schwarz need to anticipate how technology in their field
will evolve over the next decade, so they can decide where to invest. Doing this by hand
(``scenario planning'') is slow and only explores a handful of futures. Our project uses AI to
do it automatically and at much larger scale: read a body of documents, and produce a
structured map of possible futures.

\section{Where we started this sprint}
A working prototype existed for one test topic (radio-spectrum monitoring). A note from the
previous sprint flagged an open concern: the generated futures didn't seem to fall into clear,
distinct groups. The task was simply: \textbf{make it better.}

\section{What we discovered}
When we looked closely, the ``futures'' the AI produced were essentially
\textbf{indistinguishable from random} --- like shuffling the same deck of cards and calling
each shuffle a new strategy.

The root cause is a known weakness of AI language models: they are \textbf{relentlessly
positive}. When asked ``can technology A and technology B work together?'', the AI almost always
says ``yes'', because with enough effort almost anything \emph{can} be combined. But real
strategy is about \textbf{trade-offs} --- what genuinely conflicts. By never flagging conflicts,
the AI produced futures with no real backbone.

To make this concrete and measurable (rather than a matter of opinion), we built a small
\textbf{litmus test} that compares the AI's output against pure randomness. This turned a vague
worry into a hard, repeatable number --- and gave us an objective way to tell whether any later
improvement actually worked.

\section{What we did}
\textbf{(a) Quality fix --- getting honest judgments out of the AI.}
We changed \emph{how} we ask the AI to weigh technology combinations: instead of ``\emph{can}
these coexist?'' (almost always ``yes''), we now ask ``\emph{would a single, coherent design
really commit to both} --- or do they pull in opposite directions?'' That reframing makes the AI
surface the real trade-offs it was glossing over. The litmus test confirmed the improvement: the
output went from ``indistinguishable from random'' to carrying \textbf{real, measurable
structure}, and the result held up across repeated runs.

\textbf{(b) The bigger shift --- from one-off prototype to a reusable tool.}
Mid-sprint we confirmed with the product owner what actually matters to him: he is \textbf{not}
interested in the spectrum-monitoring results themselves --- that was only a test case. What he
wants is a \textbf{scalable framework}: feed it any collection of documents (any industry, any
topic) and get strategic foresight out, \textbf{without anyone having to reprogram it}. This was
effectively a scope clarification, and it reshaped our priorities for the rest of the sprint.

The problem: the tool had the demo topic (``spectrum monitoring'', specific company names, a
fixed time horizon) \textbf{baked into roughly 15 places} in its instructions. Plug in a
different topic and you'd get nonsense. So we rebuilt it the right way round. The tool now
\textbf{reads the documents first, figures out what the topic is on its own}, and feeds that
understanding into every downstream step --- like turning a single-recipe machine into one that
can read and follow \emph{any} recipe. We also added an automatic safeguard that stops anyone
from accidentally hard-coding a topic back in.

\section{The proof}
We docked a brand-new, unrelated body of documents --- \textbf{autonomous farming / agriculture
robots} --- and ran the tool. With \textbf{no code or wording changes at all}, it correctly
recognised the new field and produced a sensible analysis (sensing crops, navigating fields,
deciding where to spray, and so on). That is the core promise of the project --- ``scaling''
foresight across topics --- demonstrated end to end.

\section{Status \& quality}
\begin{itemize}
  \item \textbf{Done and working:} the topic-detection engine, the quality fix, the full
        rebuild, and the cross-topic proof.
  \item \textbf{Quality:} 109 automated tests pass, including the safeguard test above.
  \item \textbf{Not yet finalised:} the work is complete but \textbf{not yet merged} into the
        shared codebase --- a small, deliberate next step pending a quick review.
  \item \textbf{Honest scope note:} we proved the \emph{analysis core} on the second topic end to
        end. The final ``write-up'' steps (polished narrative text and scoring) use the exact same
        mechanism and are tested, but were not separately re-run on the farming example --- and
        that polished text is the part the product owner has said matters least.
\end{itemize}

\section{Next steps / open decisions}
\begin{enumerate}
  \item \textbf{Merge the work} into the shared codebase (small review, then done).
  \item \textbf{Optional:} update the dashboard/UI, which still shows labels from the old demo topic.
  \item \textbf{Optional:} run the full end-to-end output on a second topic for a complete showcase.
\end{enumerate}

\section{Mini-glossary}
\begin{itemize}
  \item \textbf{Scenario / future:} one plausible way the technology landscape could look in about a decade.
  \item \textbf{Framework:} a reusable tool that works for many topics, not a one-off built for a single one.
  \item \textbf{Knowledge base (KB):} the collection of documents we feed in as the tool's source material.
  \item \textbf{Trade-off:} a real conflict where choosing one option costs you another --- the
        substance that makes a strategy meaningful.
\end{itemize}

\end{document}
